\chapter{How Supervised Learning Works}
\label{ch:supervised}

\chapterguide%
  {Stages 1 and 2: prepared data, basic statistics, vectors, and distance measures.}
  {frame regression/classification, interpret losses, distinguish model families, and establish baselines.}

Supervised learning uses labeled examples to learn $f:\bm{x}\mapsto y$ for unseen data. First establish the objective, evaluation, and selection procedure; then compare model families.

\begin{intuition}
Learning with an answer key means using solved examples --- house prices or spam labels --- to infer a rule for new cases.
\end{intuition}

\section{Two kinds of labels}

\begin{description}[leftmargin=1.25in,style=nextline]
  \item[Regression] Predict a continuous value, such as price or temperature; measure error magnitude.
  \item[Classification] Predict a category, such as spam/not spam, optionally with class probabilities.
\end{description}

The label defines the task: exact temperature is regression; ``hot'' versus ``cold'' is classification. Counts and ordered categories may require specialized treatment.

\section{The supervised-learning objective}

A training dataset can be written as
\[
  \mathcal{D}={\left\{(\bm{x}_i,y_i)\right\}}_{i=1}^{n}.
\]
The model has parameters $\bm{\theta}$ and produces predictions $\hat{y}_i=f_{\bm{\theta}}(\bm{x}_i)$. Training usually minimizes the average loss over the training set,
\[
  \hat{R}(\bm{\theta})
  =\frac{1}{n}\sum_{i=1}^{n}
  \mathcal{L}\!\left(y_i,f_{\bm{\theta}}(\bm{x}_i)\right),
\]
The \term{empirical risk} averages per-example losses; training minimizes it over $\bm{\theta}$.

\begin{keyidea}
Training loss measures fit to known examples; validation and test metrics estimate performance on new ones.
\end{keyidea}

\section{Loss functions: how wrongness is measured}

A \term{loss function} assigns a numerical penalty to prediction errors.

For regression, with error $e=y-\hat{y}$, the two workhorses are
\[
  \mathcal{L}_{\text{squared}}(e)=e^2,
  \qquad
  \mathcal{L}_{\text{absolute}}(e)=|e|.
\]

Classification's \term{0--1 loss} counts errors but offers no gradient almost everywhere. Tractable surrogates include \term{cross-entropy} (\chapref{ch:classification}) and \term{hinge loss} (\chapref{ch:svm}); evaluate with task-relevant metrics.

\begin{mlfigure}
\begin{tikzpicture}[baseline=(current bounding box.north)]
\begin{axis}[
  width=0.45\linewidth, height=4.8cm, axis lines=left,
  title style={font=\small\bfseries\color{MLNavy},yshift=4pt}, title={Regression losses},
  xlabel={Error $(y-\hat y)$}, ylabel={Penalty},
  label style={font=\scriptsize\color{MLMuted}}, tick label style={font=\scriptsize\color{MLMuted}},
  legend style={font=\scriptsize,draw=none,fill=white,fill opacity=0.88,
                text opacity=1,rounded corners=1pt,inner sep=2pt,
                at={(0.5,0.94)},anchor=north,legend columns=3,
                /tikz/every even column/.append style={column sep=0.7em}},
  domain=-3:3, samples=101, xmin=-3, xmax=3, ymin=0, ymax=6.5,
  ytick={0,2,4,6}
]
\addplot[MLBlue,very thick]{x^2};              \addlegendentry{squared}
\addplot[MLTeal,very thick]{abs(x)};           \addlegendentry{absolute}
\addplot[MLOrange,very thick,dashed]{abs(x)<1 ? 0.5*x^2 : abs(x)-0.5}; \addlegendentry{Huber}
\end{axis}
\end{tikzpicture}
\hfill
\begin{tikzpicture}[baseline=(current bounding box.north)]
\begin{axis}[
  width=0.45\linewidth, height=4.8cm, axis lines=left,
  title style={font=\small\bfseries\color{MLNavy},yshift=4pt}, title={Cross-entropy loss},
  xlabel={Predicted probability of the true class}, ylabel={Penalty},
  label style={font=\scriptsize\color{MLMuted}}, tick label style={font=\scriptsize\color{MLMuted}},
  domain=0.02:1, samples=101, xmin=0, xmax=1, ymin=0, ymax=4
]
\addplot[MLRed,very thick]{-ln(x)};
\node[font=\scriptsize\color{MLRed},align=center,fill=white,fill opacity=0.88,
      text opacity=1,rounded corners=0.5mm,inner sep=1.5pt]
  at (axis cs:0.43,3.2) {confident\\and wrong};
\node[font=\scriptsize\color{MLGreen},align=center,fill=white,fill opacity=0.88,
      text opacity=1,rounded corners=0.5mm,inner sep=1.5pt]
  at (axis cs:0.78,0.75) {confident\\and right};
\end{axis}
\end{tikzpicture}
\end{mlfigure}

\begin{intuition}
Doubling an error quadruples squared loss, making large errors influential. Absolute loss grows linearly; Huber is quadratic near zero and linear farther away.

Cross-entropy grows without bound as the probability assigned to the true class approaches zero, strongly penalizing confident errors.
\end{intuition}

\begin{keyidea}
Use squared loss when large errors deserve extra weight; absolute or Huber loss reduces their influence.
\end{keyidea}

\section{Parametric and non-parametric models}

\begin{mltable}
\begin{tabularx}{\linewidth}{@{}>{\bfseries\leavevmode\color{MLNavy}}p{0.21\linewidth}p{0.38\linewidth}X@{}}
\toprule
\rowcolor{MLPaleBlue!92}
Model family & Structure and examples & Trade-off \\
\midrule
Parametric models & Fixed form and parameter count, independent of sample size: linear regression, logistic regression, naive Bayes. & Often fast and interpretable, but limited by the assumed form. \\
Non-parametric models & Complexity can grow with data: nearest neighbors stores training examples; trees add branches. & Flexible relationships, but often need more data and careful overfitting control. \\
\bottomrule
\end{tabularx}
\end{mltable}

The right parametric form can work well with little data; flexible non-parametric models often need more.

\section{Always start with a baseline}
\label{sec:baselines}

A \term{baseline} is a simple prediction rule evaluated with the same splits and metrics as other models:

\begin{itemize}
  \item for regression, predict the training-set mean (or median) of the label for every observation;
  \item for classification, predict the most common class for every observation;
  \item when domain rules exist (``flag any transaction over \$10{,}000''), evaluate those too.
\end{itemize}

Baselines reveal class imbalance and the value of added complexity. Lab~1's Titanic majority baseline reaches 62\% accuracy without fitting a model.

\begin{workedexample}
If 950 of 1,000 emails are legitimate, predicting ``not spam'' gives 95\% accuracy and zero spam recall. Report rare-class precision and recall (\chapref{ch:evaluation}).
\end{workedexample}

\begin{mlfigure}
\begin{tikzpicture}[
  node distance=0.4cm and 0.6cm,
  bx/.style={draw=MLBlue,fill=MLPaleBlue,rounded corners=1.5mm,minimum width=3.1cm,minimum height=0.95cm,align=center,font=\scriptsize\bfseries\color{MLNavy}},
  arr/.style={-{Stealth[length=2mm]},thick,draw=MLTeal}
]
\node[bx,fill=MLPaleGray,draw=MLMuted] (b0) {Baseline\\always guess the average};
\node[bx,right=of b0] (b1) {Simple model\\linear or a shallow tree};
\node[bx,right=of b1] (b2) {Stronger model\\forest or boosting};
\node[bx,right=of b2] (b3) {Tuned model\\cross-validated settings};
\draw[arr] (b0)--(b1); \draw[arr] (b1)--(b2); \draw[arr] (b2)--(b3);
\node[font=\scriptsize\itshape\color{MLMuted},below=0.05cm of b1] {stop here if it is good enough};
\end{tikzpicture}
\end{mlfigure}

\begin{analogy}
Increase complexity only when validation gains justify it. Stop at the simplest model that meets the task's needs.
\end{analogy}

\begin{pitfall}
Unexpectedly strong results warrant leakage checks: outcome-revealing identifiers, late-recorded features, or duplicates across splits.
\end{pitfall}

\section{A map of the algorithms ahead}

The next six model chapters apply the workflow from \chapref{ch:training}, organized by their core ideas.

\begin{mltable}
\renewcommand{\arraystretch}{1.16}
\setlength{\tabcolsep}{5pt}
\begin{tabularx}{\linewidth}{@{}>{\footnotesize\raggedright\arraybackslash}p{0.12\linewidth}>{\footnotesize\raggedright\arraybackslash}p{0.42\linewidth}>{\footnotesize\raggedright\arraybackslash}X@{}}
\toprule
\rowcolor{MLPaleBlue!92}
\normalfont\bfseries Chapter & \bfseries Algorithms & \bfseries Core idea \\
\midrule
\ref{ch:regression} & Linear, polynomial, ridge, lasso, elastic net regression & Fit a weighted sum of features; penalize large weights \\
\ref{ch:classification} & Logistic regression, softmax, LDA, QDA, perceptron & Draw a linear boundary between classes \\
\ref{ch:knn-nb} & $k$-nearest neighbors, naive Bayes & Predict from similar examples; predict from per-feature probabilities \\
\ref{ch:decision-trees} & Decision trees (CART) & Ask a sequence of yes/no questions \\
\ref{ch:ensembles} & Bagging, random forests, AdaBoost, gradient boosting, stacking & Combine many models into one predictor \\
\ref{ch:svm} & Support vector machines, kernels, SVR & Maximize the margin; separate after mapping to a richer space \\
\bottomrule
\end{tabularx}
\end{mltable}

\begin{keyidea}
Supervised models differ in representable functions, assumptions, and data needs. Learn these trade-offs rather than formulas alone.
\end{keyidea}

\begin{modelcard}{Supervised Learning - Project Checklist}
\begin{tabularx}{\linewidth}{@{}>{\bfseries\leavevmode\color{MLNavy}}p{0.17\linewidth}>{\raggedright\arraybackslash}X@{}}
Question 1 & Is the label a number, a category, or a special type such as a count or ordered category? \cr
Question 2 & What simple baseline should the model beat? \cr
Question 3 & Which loss guides training, and which metric matches the real cost of mistakes? \cr
Question 4 & How will the split imitate future predictions? \cr
Question 5 & Does the final model need interpretability, probabilities, speed, or a nonlinear boundary? \cr
\end{tabularx}
\end{modelcard}

\begin{quickcheck}
\begin{enumerate}
  \item How does a regression label differ from a classification label, and what does that change about the kind of prediction produced?
  \item Why can the loss used for fitting differ from the metric used to report practical performance?
  \item What does a learned model have to demonstrate beyond beating its own training loss?
\end{enumerate}
\end{quickcheck}
